\documentclass[12pt]{article}
\usepackage{fullpage,amsmath,amssymb,amsthm,hyperref,amsfonts}
\newtheorem{theorem}{Theorem}
\newtheorem{lemma}[theorem]{Lemma}
\newtheorem{proposition}[theorem]{Proposition}
\theoremstyle{definition}
\newtheorem{definition}[theorem]{Definition}
\newtheorem{remark}[theorem]{Remark}

\title{Problem 9: Algebraic Relations on Determinantal Tensors}
\author{}
\date{}

\begin{document}
\maketitle

\section{Statement and Answer}

Fix $n\geq 5$. Let $A^{(1)},\ldots,A^{(n)}\in\mathbb{R}^{3\times 4}$ be Zariski-generic. For
$\alpha,\beta,\gamma,\delta\in[n]:=\{1,\ldots,n\}$ and $i,j,k,\ell\in\{1,2,3\}$, define
\begin{equation}\label{eq:Q-def}
Q^{(\alpha\beta\gamma\delta)}_{ijk\ell}\;=\;
\det\!\begin{bmatrix}
A^{(\alpha)}(i,:)\\
A^{(\beta)}(j,:)\\
A^{(\gamma)}(k,:)\\
A^{(\delta)}(\ell,:)
\end{bmatrix},
\end{equation}
where $A(i,:)\in\mathbb{R}^4$ is the $i$-th row of $A$ and the bracket denotes vertical
concatenation. Given a scalar array $\lambda\in\mathbb{R}^{n\times n\times n\times n}$ with
$\lambda_{\alpha\beta\gamma\delta}\neq 0$ exactly when $(\alpha,\beta,\gamma,\delta)$ are not
all equal, set
\begin{equation}\label{eq:Pdef}
P^{(\alpha\beta\gamma\delta)}_{ijk\ell}\;=\;\lambda_{\alpha\beta\gamma\delta}\,
Q^{(\alpha\beta\gamma\delta)}_{ijk\ell}.
\end{equation}

We are asked: does there exist a polynomial map
$\mathbf{F}:\mathbb{R}^{81n^4}\to\mathbb{R}^N$ with
\begin{enumerate}
\item[(P1)] $\mathbf{F}$ is independent of $A^{(1)},\ldots,A^{(n)}$;
\item[(P2)] the degrees of the coordinate functions of $\mathbf{F}$ are independent of $n$;
\item[(P3)] $\mathbf{F}\!\left(\{P^{(\alpha\beta\gamma\delta)}\}_{\alpha,\beta,\gamma,\delta}\right)=0$
if and only if there exist $u,v,w,x\in(\mathbb{R}^*)^n$ with
$\lambda_{\alpha\beta\gamma\delta}=u_\alpha v_\beta w_\gamma x_\delta$ for all
$(\alpha,\beta,\gamma,\delta)$ that are not all equal?
\end{enumerate}

\textbf{Answer.} \emph{Yes.} An explicit construction is given in Section~\ref{sec:construction}.

\section{Construction}\label{sec:construction}

The big array $\mathbf{P}=\{P^{(\alpha\beta\gamma\delta)}_{ijk\ell}\}$ from \eqref{eq:Pdef} is
naturally indexed by four pairs $(\alpha,i),(\beta,j),(\gamma,k),(\delta,\ell)$ with
$\alpha,\beta,\gamma,\delta\in[n]$ and $i,j,k,\ell\in[3]$. We regard $\mathbf{P}$ as a
$4$-tensor of size $3n\times 3n\times 3n\times 3n$, with mode-$1$ index $I_1=(\alpha,i)$,
mode-$2$ index $I_2=(\beta,j)$, mode-$3$ index $I_3=(\gamma,k)$, mode-$4$ index
$I_4=(\delta,\ell)$.

\begin{definition}[Mode-$s$ unfolding]\label{def:unfold}
For each $s\in\{1,2,3,4\}$, the \emph{mode-$s$ unfolding} $\mathbf{M}^{(s)}\in\mathbb{R}^{3n\times(3n)^3}$
of $\mathbf{P}$ is the matrix whose $(I_s,(I_{s_1},I_{s_2},I_{s_3}))$-entry equals
$\mathbf{P}_{I_1 I_2 I_3 I_4}$, where $(s_1,s_2,s_3)$ is the increasing enumeration of
$\{1,2,3,4\}\setminus\{s\}$ and the column index $(I_{s_1},I_{s_2},I_{s_3})$ runs over
$[3n]^3$ in lexicographic order.
\end{definition}

\begin{definition}[Construction of $\mathbf{F}$]\label{def:F}
Let $\mathbf{F}:\mathbb{R}^{81n^4}\to\mathbb{R}^N$ be the polynomial map whose coordinate
functions are all $5\times 5$ minors of $\mathbf{M}^{(s)}$, for every $s\in\{1,2,3,4\}$.
Equivalently, the coordinates of $\mathbf{F}$ are
\[
\bigl\{\det \mathbf{M}^{(s)}[R\,|\,C]\;:\;s\in\{1,2,3,4\},\;R\in\binom{[3n]}{5},\;C\in\binom{[3n]^3}{5}\bigr\}.
\]
The total number of coordinates is
$N=4\cdot\binom{3n}{5}\cdot\binom{(3n)^3}{5}$.
\end{definition}

Each $5\times 5$ minor is a polynomial of degree $5$ in the $81n^4$ entries of the input
array $\mathbf{P}$, with coefficients $\pm 1$ (independent of $A^{(\alpha)}$ and of $n$).
Properties (P1) and (P2) are immediate. The remainder of the paper proves (P3).

\section{Generic non-degeneracy of the matrices $A^{(\alpha)}$}\label{sec:generic}

The hypothesis ``Zariski-generic'' means that $(A^{(1)},\ldots,A^{(n)})$ lies outside a
proper Zariski-closed subset of $\mathbb{R}^{12n}$. We fix once and for all a finite list
of polynomial conditions that hold throughout the proof; on the complement of their
common zero locus we assert all conclusions.

\begin{lemma}\label{lem:generic-AB}
For Zariski-generic $A^{(1)},\ldots,A^{(n)}\in\mathbb{R}^{3\times 4}$, the following hold.
\begin{enumerate}
\item[(G1)] For every $\alpha\in[n]$, the three rows of $A^{(\alpha)}$ are
$\mathbb{R}$-linearly independent (so $A^{(\alpha)}$ has rank $3$).
\item[(G2)] For every distinct $\alpha,\beta\in[n]$ and every choice
$i_0,j_0\in\{1,2,3\}$, the row vectors $A^{(\alpha)}(i_0,:),A^{(\beta)}(j_0,:)\in\mathbb{R}^4$
are linearly independent.
\item[(G3)] For every choice of three distinct indices $\alpha,\beta,\gamma\in[n]$ and
every $i_0,j_0,k_0\in\{1,2,3\}$, the rows
$A^{(\alpha)}(i_0,:),A^{(\beta)}(j_0,:),A^{(\gamma)}(k_0,:)$ are linearly independent.
\item[(G4)] For every choice of four distinct indices $\alpha,\beta,\gamma,\delta\in[n]$
and every $i_0,j_0,k_0,\ell_0\in\{1,2,3\}$ we have
\[
Q^{(\alpha\beta\gamma\delta)}_{i_0 j_0 k_0 \ell_0}\neq 0.
\]
\end{enumerate}
\end{lemma}

\begin{proof}
Each of (G1)--(G4) is the non-vanishing of a single polynomial in the entries of
$A^{(1)},\ldots,A^{(n)}$ (a determinant of a $3\times 3$, $2\times 4$ minor sub-determinant,
$3\times 4$ minor sub-determinant, or $4\times 4$ matrix, respectively, all of which are
not identically zero as polynomials in the entries since one can specialize the rows to a
basis of $\mathbb{R}^4$ to make them nonzero). Finitely many such conditions intersect in a
non-empty Zariski-open subset of $\mathbb{R}^{12n}$. The hypothesis of generic
$A^{(\alpha)}$ ensures the chosen tuple lies in this open set.
\end{proof}

\begin{lemma}\label{lem:generic-cofactor}
For Zariski-generic $A^{(1)},\ldots,A^{(n)}\in\mathbb{R}^{3\times 4}$, additionally:
\begin{enumerate}
\item[(G5)] For every choice of three distinct $\beta,\gamma,\delta\in[n]$ and every
$j_0,k_0,\ell_0\in[3]$, the $3\times 3$ minors
$M^{(p)}=M^{(p)}(\beta,\gamma,\delta,j_0,k_0,\ell_0)$
($p=1,2,3,4$) of the $3\times 4$ matrix
\[
\begin{bmatrix}A^{(\beta)}(j_0,:)\\ A^{(\gamma)}(k_0,:)\\ A^{(\delta)}(\ell_0,:)\end{bmatrix}\in\mathbb{R}^{3\times 4}
\]
obtained by deleting column $p$ are not all zero, and at least two of them are nonzero.
\end{enumerate}
\end{lemma}

\begin{proof}
By (G3), the three rows $A^{(\beta)}(j_0,:),A^{(\gamma)}(k_0,:),A^{(\delta)}(\ell_0,:)\in\mathbb{R}^4$
are linearly independent for distinct $\beta,\gamma,\delta$, so the corresponding $3\times 4$
matrix has rank $3$. Hence at least one $3\times 3$ minor is nonzero. The condition that
``at most one $M^{(p)}$ is nonzero'' is a closed condition (vanishing of all but one
specified minor), and it does not hold identically on $\mathbb{R}^{3\times 4}$ matrices of
rank $3$ (the matrix $\begin{bmatrix}1&0&0&0\\0&1&0&0\\0&0&1&0\end{bmatrix}$ shows three
minors are zero and one nonzero, while $\begin{bmatrix}1&0&0&0\\0&1&0&0\\0&0&1&1\end{bmatrix}$
makes two minors nonzero). So a Zariski-open further restriction enforces (G5).
\end{proof}

Throughout the rest of the proof, we assume $(A^{(1)},\ldots,A^{(n)})$ satisfies
(G1)--(G5). We refer to this collectively as the \emph{genericity hypothesis}.

\section{Proof of (P1) and (P2)}\label{sec:P1P2}

\begin{proposition}[Properties (P1) and (P2)]\label{prop:P1P2}
The map $\mathbf{F}$ from Definition~\ref{def:F} satisfies (P1) and (P2).
\end{proposition}

\begin{proof}
Each coordinate of $\mathbf{F}$ is a $5\times 5$ minor of an unfolding $\mathbf{M}^{(s)}$, whose
entries are precisely the entries of $\mathbf{P}$. Thus each coordinate is a polynomial of
total degree exactly $5$ in the $81n^4$ entries $P^{(\alpha\beta\gamma\delta)}_{ijk\ell}$,
with coefficients $\pm 1$. The polynomials do not refer to $A^{(\alpha)}$ at all, proving
(P1). The degree is $5$ regardless of $n$, proving (P2).
\end{proof}

\section{The forward direction of (P3)}\label{sec:P3-fwd}

In this section we prove the implication
\[
\bigl(\exists\, u,v,w,x\in(\mathbb{R}^*)^n\colon \lambda_{\alpha\beta\gamma\delta}=u_\alpha v_\beta w_\gamma x_\delta\bigr)
\;\Longrightarrow\;\mathbf{F}(\mathbf{P})=0.
\]
The strategy is to show that the rank of the mode-$1$ unfolding $\mathbf{M}^{(1)}$ is at
most $4$, hence all its $5\times 5$ minors vanish; symmetry handles modes $2,3,4$.

\begin{lemma}[Rank-$4$ bound under multiplicative $\lambda$]\label{lem:fwd-rank}
Suppose $\lambda_{\alpha\beta\gamma\delta}=u_\alpha v_\beta w_\gamma x_\delta$ for some
$u,v,w,x\in(\mathbb{R}^*)^n$ and all $(\alpha,\beta,\gamma,\delta)$ that are not all equal.
Set $\lambda_{\alpha\alpha\alpha\alpha}=u_\alpha v_\alpha w_\alpha x_\alpha$ as well (this is
allowed because $Q^{(\alpha\alpha\alpha\alpha)}_{ijk\ell}=0$ identically; see
Lemma~\ref{lem:Q-skew}). Then
\[
\mathrm{rank}\bigl(\mathbf{M}^{(s)}\bigr)\leq 4\qquad\text{for every }s\in\{1,2,3,4\}.
\]
\end{lemma}

We first record the basic vanishing.

\begin{lemma}[Vanishing of $Q$ on the diagonal]\label{lem:Q-skew}
For every $\alpha\in[n]$ and every $i,j,k,\ell\in\{1,2,3\}$,
$Q^{(\alpha\alpha\alpha\alpha)}_{ijk\ell}=0$.
\end{lemma}

\begin{proof}
The matrix $\bigl[A^{(\alpha)}(i,:);A^{(\alpha)}(j,:);A^{(\alpha)}(k,:);A^{(\alpha)}(\ell,:)\bigr]$
has $4$ rows, each from the $3$-row matrix $A^{(\alpha)}\in\mathbb{R}^{3\times 4}$. By the
pigeonhole principle two of $i,j,k,\ell$ coincide, so two rows of the $4\times 4$ matrix are
equal, and its determinant is $0$.
\end{proof}

\begin{proof}[Proof of Lemma~\ref{lem:fwd-rank}]
Fix $s=1$ (the case $s\in\{2,3,4\}$ is identical by relabeling modes). Define rescaled rows
\begin{equation}\label{eq:rescaled-rows}
\widetilde{a}^{(\alpha)}_i\;:=\;u_\alpha\, A^{(\alpha)}(i,:),\quad
\widetilde{b}^{(\beta)}_j\;:=\;v_\beta\, A^{(\beta)}(j,:),\quad
\widetilde{c}^{(\gamma)}_k\;:=\;w_\gamma\, A^{(\gamma)}(k,:),\quad
\widetilde{d}^{(\delta)}_\ell\;:=\;x_\delta\, A^{(\delta)}(\ell,:)
\end{equation}
in $\mathbb{R}^4$. By multilinearity of the determinant in its rows,
\begin{align}
P^{(\alpha\beta\gamma\delta)}_{ijk\ell}
&=u_\alpha v_\beta w_\gamma x_\delta\det\!\begin{bmatrix}A^{(\alpha)}(i,:)\\A^{(\beta)}(j,:)\\A^{(\gamma)}(k,:)\\A^{(\delta)}(\ell,:)\end{bmatrix}\notag\\
&=\det\!\begin{bmatrix}\widetilde{a}^{(\alpha)}_i\\ \widetilde{b}^{(\beta)}_j\\ \widetilde{c}^{(\gamma)}_k\\ \widetilde{d}^{(\delta)}_\ell\end{bmatrix}.\label{eq:multilin}
\end{align}
Equation~\eqref{eq:multilin} also holds when $\alpha=\beta=\gamma=\delta$ by
Lemma~\ref{lem:Q-skew} and the corresponding repetition of rows in the rescaled matrix.

Expand the determinant on the right of \eqref{eq:multilin} along the first row using cofactor
expansion. For $p\in\{1,2,3,4\}$, let
\begin{equation}\label{eq:Mp}
M^{(p)}_{(\beta,j),(\gamma,k),(\delta,\ell)}
\;:=\;\det\!\begin{bmatrix}
\widetilde{b}^{(\beta)}_j[\widehat{p}]\\
\widetilde{c}^{(\gamma)}_k[\widehat{p}]\\
\widetilde{d}^{(\delta)}_\ell[\widehat{p}]
\end{bmatrix},
\end{equation}
where $v[\widehat{p}]\in\mathbb{R}^3$ denotes the vector $v\in\mathbb{R}^4$ with its $p$-th
entry deleted. Then by Laplace expansion along the top row:
\begin{equation}\label{eq:Laplace}
P^{(\alpha\beta\gamma\delta)}_{ijk\ell}
=\sum_{p=1}^{4}(-1)^{p+1}\bigl(\widetilde{a}^{(\alpha)}_i\bigr)_p\cdot
M^{(p)}_{(\beta,j),(\gamma,k),(\delta,\ell)}.
\end{equation}

Now read \eqref{eq:Laplace} as a statement about rows of $\mathbf{M}^{(1)}$. The row of
$\mathbf{M}^{(1)}$ indexed by $(\alpha,i)\in[n]\times[3]$ has entries indexed by
$\bigl((\beta,j),(\gamma,k),(\delta,\ell)\bigr)\in[n]^3\times[3]^3$, namely
$P^{(\alpha\beta\gamma\delta)}_{ijk\ell}$. Define a $4\times(3n)^3$ matrix $\mathbf{N}$ whose
$p$-th row has $\bigl((\beta,j),(\gamma,k),(\delta,\ell)\bigr)$-entry equal to
$(-1)^{p+1}M^{(p)}_{(\beta,j),(\gamma,k),(\delta,\ell)}$. Note that $\mathbf{N}$ does not
depend on $\alpha$ or $i$. Equation~\eqref{eq:Laplace} reads, for each fixed
$(\alpha,i)$ and column index $C=\bigl((\beta,j),(\gamma,k),(\delta,\ell)\bigr)$,
\begin{equation}\label{eq:row-decomp}
\mathbf{M}^{(1)}\bigl[(\alpha,i),C\bigr]
=\sum_{p=1}^{4}\bigl(\widetilde{a}^{(\alpha)}_i\bigr)_p\cdot \mathbf{N}[p,C]
=\bigl(\widetilde{a}^{(\alpha)}_i\bigr)\cdot\mathbf{N}[:,C],
\end{equation}
where $\widetilde{a}^{(\alpha)}_i\in\mathbb{R}^{1\times 4}$ is regarded as a row vector.
Hence in matrix form
\begin{equation}\label{eq:matrix-factor}
\mathbf{M}^{(1)}\;=\;\widetilde{\mathbf{A}}\,\mathbf{N},
\qquad\text{where }\widetilde{\mathbf{A}}\in\mathbb{R}^{3n\times 4}\text{ has rows }\widetilde{a}^{(\alpha)}_i.
\end{equation}
The matrix $\widetilde{\mathbf{A}}$ has only $4$ columns; therefore
$\mathrm{rank}(\widetilde{\mathbf{A}})\leq 4$ and consequently
\begin{equation}\label{eq:M1-rank}
\mathrm{rank}\bigl(\mathbf{M}^{(1)}\bigr)\leq\min\bigl(\mathrm{rank}(\widetilde{\mathbf{A}}),\mathrm{rank}(\mathbf{N})\bigr)\leq 4.
\end{equation}

For $s\in\{2,3,4\}$, the same argument applies with the roles of the modes permuted:
expand the determinant along the row associated with mode $s$ to obtain
$\mathbf{M}^{(s)}=\widetilde{\mathbf{B}}_s\,\mathbf{N}_s$ with $\widetilde{\mathbf{B}}_s$ of
size $3n\times 4$. Hence $\mathrm{rank}(\mathbf{M}^{(s)})\leq 4$ for every $s$.
\end{proof}

\begin{proposition}[Forward direction of (P3)]\label{prop:fwd}
If $\lambda_{\alpha\beta\gamma\delta}=u_\alpha v_\beta w_\gamma x_\delta$ for some
$u,v,w,x\in(\mathbb{R}^*)^n$ and all non-all-equal $(\alpha,\beta,\gamma,\delta)$, then
$\mathbf{F}(\mathbf{P})=0$.
\end{proposition}

\begin{proof}
By Lemma~\ref{lem:fwd-rank}, $\mathrm{rank}(\mathbf{M}^{(s)})\leq 4$ for every $s\in\{1,2,3,4\}$.
A matrix of rank $\leq 4$ has every $5\times 5$ minor equal to $0$. Since the coordinates
of $\mathbf{F}$ are exactly these $5\times 5$ minors over $s=1,2,3,4$, we have
$\mathbf{F}(\mathbf{P})=0$.
\end{proof}

\section{The reverse direction of (P3)}\label{sec:P3-rev}

We now prove the converse: if $\mathbf{F}(\mathbf{P})=0$ (i.e., every mode-$s$ unfolding has
rank $\leq 4$) and the genericity hypothesis on $A^{(\alpha)}$ holds, then there exist
$u,v,w,x\in(\mathbb{R}^*)^n$ such that $\lambda_{\alpha\beta\gamma\delta}=u_\alpha v_\beta w_\gamma x_\delta$
for every non-all-equal $(\alpha,\beta,\gamma,\delta)$.

\subsection{Reducing rank-$4$ to a structural identity on $\lambda$}

\begin{lemma}[Mode-$1$ row decomposition for arbitrary $\lambda$]\label{lem:abs-rowdecomp}
With the notation of Section~\ref{sec:construction}, for every $\alpha,\beta,\gamma,\delta\in[n]$
and $i,j,k,\ell\in\{1,2,3\}$,
\begin{equation}\label{eq:abs-Laplace}
P^{(\alpha\beta\gamma\delta)}_{ijk\ell}
=\lambda_{\alpha\beta\gamma\delta}\sum_{p=1}^{4}(-1)^{p+1}A^{(\alpha)}(i,p)
\cdot M^{(p)}_{(\beta,j),(\gamma,k),(\delta,\ell)},
\end{equation}
where now $M^{(p)}_{(\beta,j),(\gamma,k),(\delta,\ell)}$ denotes the $3\times 3$ determinant
of the matrix obtained by deleting column $p$ from
$\bigl[A^{(\beta)}(j,:);A^{(\gamma)}(k,:);A^{(\delta)}(\ell,:)\bigr]\in\mathbb{R}^{3\times 4}$.
\end{lemma}

\begin{proof}
Expand $Q^{(\alpha\beta\gamma\delta)}_{ijk\ell}=\det[A^{(\alpha)}(i,:);A^{(\beta)}(j,:);A^{(\gamma)}(k,:);A^{(\delta)}(\ell,:)]$
along the first row by Laplace's formula and multiply by $\lambda_{\alpha\beta\gamma\delta}$.
\end{proof}

\begin{lemma}[Image of mode-$1$ rows]\label{lem:row-image}
Fix any $\alpha\in[n]$ and consider the rows $\{R^{(\alpha,i)}\}_{i=1}^{3}$ of
$\mathbf{M}^{(1)}$ indexed by $(\alpha,1),(\alpha,2),(\alpha,3)$. For every column index
$C=\bigl((\beta,j),(\gamma,k),(\delta,\ell)\bigr)$,
\begin{equation}\label{eq:row-formula}
R^{(\alpha,i)}[C]=\lambda_{\alpha\beta\gamma\delta}\,\bigl\langle A^{(\alpha)}(i,:),\,m(C)\bigr\rangle,
\end{equation}
where $m(C)\in\mathbb{R}^4$ is the vector with $p$-th coordinate
$(-1)^{p+1}M^{(p)}_{(\beta,j),(\gamma,k),(\delta,\ell)}$ from
Lemma~\ref{lem:abs-rowdecomp}.
\end{lemma}

\begin{proof}
This is exactly the statement of Lemma~\ref{lem:abs-rowdecomp} written using the inner
product on $\mathbb{R}^4$.
\end{proof}

\subsection{Mode-$1$ rank constraint forces a multiplicative structure in $\alpha$}

We now exploit the assumption $\mathrm{rank}(\mathbf{M}^{(1)})\leq 4$.

\begin{lemma}[Three-row independence]\label{lem:three-rows}
For every $\alpha\in[n]$, the three rows
$R^{(\alpha,1)},R^{(\alpha,2)},R^{(\alpha,3)}$ of $\mathbf{M}^{(1)}$ are linearly independent
over $\mathbb{R}$ (in $\mathbb{R}^{(3n)^3}$).
\end{lemma}

\begin{proof}
Suppose $\sum_{i=1}^{3}c_i R^{(\alpha,i)}=0$ for some $c_1,c_2,c_3\in\mathbb{R}$. Choose any
distinct $\beta,\gamma,\delta\in[n]\setminus\{\alpha\}$ (possible because $n\geq 5\geq 4$)
and any $j_0,k_0,\ell_0\in\{1,2,3\}$. By (G4), $Q^{(\alpha\beta\gamma\delta)}_{i\,j_0 k_0 \ell_0}\neq 0$
for at least some $i$, but more is true: $\lambda_{\alpha\beta\gamma\delta}\neq 0$ since
$(\alpha,\beta,\gamma,\delta)$ are all distinct (hence not all equal). Reading
\eqref{eq:row-formula} at the column $C_0=\bigl((\beta,j_0),(\gamma,k_0),(\delta,\ell_0)\bigr)$,
\[
\sum_{i=1}^{3}c_i\, R^{(\alpha,i)}[C_0]
=\lambda_{\alpha\beta\gamma\delta}\,\Bigl\langle\sum_{i=1}^{3}c_i A^{(\alpha)}(i,:),\,m(C_0)\Bigr\rangle.
\]
This vanishes for every choice of the column $C_0$ as above; that is, the vector
$v:=\sum_{i=1}^{3}c_i A^{(\alpha)}(i,:)\in\mathbb{R}^4$ satisfies
$\langle v,m(C_0)\rangle=0$ for every $\beta,\gamma,\delta\in[n]\setminus\{\alpha\}$
distinct, and every $j_0,k_0,\ell_0\in\{1,2,3\}$.

We claim that the vectors $\{m(C_0)\}$, as $(\beta,\gamma,\delta)$ ranges over distinct triples
in $[n]\setminus\{\alpha\}$ and $(j_0,k_0,\ell_0)$ over $[3]^3$, span $\mathbb{R}^4$. Indeed,
fix any three distinct $\beta,\gamma,\delta\in[n]\setminus\{\alpha\}$ (this is possible
since $|[n]\setminus\{\alpha\}|=n-1\geq 4$). By (G3) the $3\times 4$ matrix
$\bigl[A^{(\beta)}(j_0,:);A^{(\gamma)}(k_0,:);A^{(\delta)}(\ell_0,:)\bigr]$ has rank $3$ for
each fixed $(j_0,k_0,\ell_0)$. Its $4$ minors $M^{(p)}$ (for $p=1,2,3,4$) form the entries
of $m(C_0)$. By Cramer's rule, $m(C_0)$ is (up to sign) a normal vector to the row span
of the $3\times 4$ matrix, hence nonzero. Now vary $(j_0,k_0,\ell_0)\in[3]^3$ for fixed
$(\beta,\gamma,\delta)$: the vectors $m(C_0)$ are normals to the $3$-dimensional row spans
of all $3\times 4$ matrices formed by choosing one row each from $A^{(\beta)},A^{(\gamma)},A^{(\delta)}$.
By varying these choices (and additionally varying $\beta,\gamma,\delta$) the resulting
normals span $\mathbb{R}^4$ generically; explicitly, the set of all such $m(C_0)$ is the
common image of polynomial maps $\mathbb{R}^{12n}\to\mathbb{R}^4$, one for each combinatorial
choice, and the open condition ``the union of these images spans $\mathbb{R}^4$'' is
non-empty (witnessed by, e.g., $A^{(\alpha)}(i,p)=\delta_{ip}$ extended by independent rows)
and so forms a Zariski-open subset of the parameter space; we add this condition to (G5).

Therefore $v$ is orthogonal to a spanning set of $\mathbb{R}^4$, hence $v=0$, i.e.,
$\sum_{i=1}^{3}c_i A^{(\alpha)}(i,:)=0$. By (G1) the rows of $A^{(\alpha)}$ are linearly
independent, so $c_1=c_2=c_3=0$.
\end{proof}

\begin{lemma}[Mode-$1$ row span]\label{lem:row-span}
Suppose $\mathrm{rank}(\mathbf{M}^{(1)})\leq 4$. Then there exists a $4$-dimensional subspace
$V_1\subseteq\mathbb{R}^{(3n)^3}$ that contains every row $R^{(\alpha,i)}$ of $\mathbf{M}^{(1)}$.
\end{lemma}

\begin{proof}
By Lemma~\ref{lem:three-rows} the rows $R^{(\alpha,1)},R^{(\alpha,2)},R^{(\alpha,3)}$ are
linearly independent for every $\alpha$. Since $\mathrm{rank}(\mathbf{M}^{(1)})\leq 4$ and
the row span is exactly the rank (over $\mathbb{R}$), the row span is a subspace of dimension
$\leq 4$. Take $V_1$ to be the row span.
\end{proof}

\begin{lemma}[Two distinct $\alpha$'s already span $V_1$]\label{lem:two-alpha-span}
Fix any two distinct $\alpha,\alpha'\in[n]$. Under the genericity hypothesis and
$\mathrm{rank}(\mathbf{M}^{(1)})\leq 4$, the six rows
$\{R^{(\alpha,i)},R^{(\alpha',i')}:i,i'\in\{1,2,3\}\}$ have linear span equal to $V_1$ if at
least one of these six rows is nonzero, and the span has dimension exactly $4$.
\end{lemma}

\begin{proof}
The two sets of three rows each have dimension $3$ by Lemma~\ref{lem:three-rows}. Their
combined span is at least $3$ and at most $\dim V_1=4$ (by Lemma~\ref{lem:row-span}). It is
strictly more than $3$: otherwise the rows for $\alpha'$ lie in the $3$-dim span of rows for
$\alpha$, hence $R^{(\alpha',i')}=\sum_{i=1}^{3}c_{i,i'}R^{(\alpha,i)}$ for scalars $c_{i,i'}$.
Evaluating both sides at column $C_0=\bigl((\beta,j_0),(\gamma,k_0),(\delta,\ell_0)\bigr)$
with distinct $\beta,\gamma,\delta\in[n]\setminus\{\alpha,\alpha'\}$ (possible since
$n\geq 5$) gives
\[
\lambda_{\alpha'\beta\gamma\delta}\,\bigl\langle A^{(\alpha')}(i',:),m(C_0)\bigr\rangle
=\sum_{i}c_{i,i'}\lambda_{\alpha\beta\gamma\delta}\,\bigl\langle A^{(\alpha)}(i,:),m(C_0)\bigr\rangle.
\]
Since the vectors $\{m(C_0)\}$ span $\mathbb{R}^4$ (proof of Lemma~\ref{lem:three-rows}),
this enforces a linear relation between the row $A^{(\alpha')}(i',:)$ and the rows
$A^{(\alpha)}(i,:)$ of the form
$\lambda_{\alpha'\beta\gamma\delta} A^{(\alpha')}(i',:)
=\sum_{i}c_{i,i'}\lambda_{\alpha\beta\gamma\delta}A^{(\alpha)}(i,:)$
for every triple $(\beta,\gamma,\delta)$ of distinct indices outside $\{\alpha,\alpha'\}$. By
(G2), $A^{(\alpha)}(i,:)$ and $A^{(\alpha')}(i',:)$ are linearly independent, contradicting
this. So the combined span has dimension at least $4$, and being $\leq 4$ it equals $V_1$.
\end{proof}

\begin{lemma}[Pointwise scalar law for $\lambda$ in the first index]\label{lem:lambda-firstindex}
Assume the genericity hypothesis and $\mathrm{rank}(\mathbf{M}^{(1)})\leq 4$. For every
$\alpha,\alpha'\in[n]$ and every $(\beta,\gamma,\delta)\in[n]^3$ with
$\{\beta,\gamma,\delta\}\cap\{\alpha,\alpha'\}=\varnothing$ and
$\beta,\gamma,\delta$ pairwise distinct, the ratio
\begin{equation}\label{eq:ratio}
r_{\alpha,\alpha'}\;:=\;\frac{\lambda_{\alpha\beta\gamma\delta}}{\lambda_{\alpha'\beta\gamma\delta}}
\end{equation}
is well-defined and independent of $(\beta,\gamma,\delta)$.
\end{lemma}

\begin{proof}
Both numerator and denominator are nonzero because $(\alpha,\beta,\gamma,\delta)$ and
$(\alpha',\beta,\gamma,\delta)$ are tuples of four distinct indices, hence not all equal,
hence $\lambda_{\cdot\,\beta\gamma\delta}\neq 0$ by hypothesis.

Fix two such triples $(\beta,\gamma,\delta)$ and $(\beta',\gamma',\delta')$ disjoint from
$\{\alpha,\alpha'\}$, both with distinct entries. We must show
\[
\frac{\lambda_{\alpha\beta\gamma\delta}}{\lambda_{\alpha'\beta\gamma\delta}}
=\frac{\lambda_{\alpha\beta'\gamma'\delta'}}{\lambda_{\alpha'\beta'\gamma'\delta'}}.
\]
By Lemma~\ref{lem:row-span} the rows $R^{(\alpha,i)}$ and $R^{(\alpha',i')}$ all lie in
$V_1$ (a $4$-dimensional space). By Lemma~\ref{lem:two-alpha-span}, the rows
$R^{(\alpha,1)},R^{(\alpha,2)},R^{(\alpha,3)}$ together with any $R^{(\alpha',i')}$ for some
$i'\in\{1,2,3\}$ span $V_1$. Hence we may fix $i'_0$ such that
$\{R^{(\alpha,1)},R^{(\alpha,2)},R^{(\alpha,3)},R^{(\alpha',i'_0)}\}$ form a basis of $V_1$.

Consequently for each $i\in\{1,2,3\}$ there exist unique scalars
$\mu^{(\alpha,\alpha')}_{i'_0}, \nu^{(\alpha,\alpha')}_{i,i'_0}\in\mathbb{R}$ such that
\begin{equation}\label{eq:basis-rep}
R^{(\alpha',i'_0)}=\sum_{i=1}^{3}\nu^{(\alpha,\alpha')}_{i,i'_0}R^{(\alpha,i)}+\mu^{(\alpha,\alpha')}_{i'_0}\cdot 0,
\end{equation}
which is a tautology. The non-tautological assertion comes from comparing the two columns
$C_0=\bigl((\beta,j_0),(\gamma,k_0),(\delta,\ell_0)\bigr)$ and $C'_0=\bigl((\beta',j_0),(\gamma',k_0),(\delta',\ell_0)\bigr)$
when both $(\alpha,\beta,\gamma,\delta)$ and $(\alpha',\beta,\gamma,\delta)$ have all four
entries distinct. We use a more direct extraction below; the basis
existence above is recorded only to confirm we are working inside the $4$-dim space.

Define for fixed $(\beta,\gamma,\delta)$ disjoint from $\{\alpha,\alpha'\}$, with all distinct
entries, the $3$-vectors $u_{\alpha,(\beta\gamma\delta)},u_{\alpha',(\beta\gamma\delta)}\in\mathbb{R}^3$
indexed by $(j,k,\ell)\in[3]^3$ as follows: take a fixed reference $(j_0,k_0,\ell_0)\in[3]^3$
with $M^{(p_0)}_{(\beta,j_0),(\gamma,k_0),(\delta,\ell_0)}\neq 0$ for some $p_0$ (possible
by (G5)). Then by Lemma~\ref{lem:abs-rowdecomp}
\begin{align}
\mathbf{M}^{(1)}\bigl[(\alpha,i),C_0\bigr]
&=\lambda_{\alpha\beta\gamma\delta}\,\sum_{p}(-1)^{p+1}A^{(\alpha)}(i,p)M^{(p)}_{(\beta,j_0),(\gamma,k_0),(\delta,\ell_0)}\notag\\
&=\lambda_{\alpha\beta\gamma\delta}\,\bigl\langle A^{(\alpha)}(i,:),m(\beta,\gamma,\delta;j_0,k_0,\ell_0)\bigr\rangle,\label{eq:M1-eval}
\end{align}
and similarly for $(\alpha',i)$ with prefactor $\lambda_{\alpha'\beta\gamma\delta}$.

Now consider the $6\times 1$ column vector
$\mathbf{M}^{(1)}[(\alpha,1),C_0],\mathbf{M}^{(1)}[(\alpha,2),C_0],\mathbf{M}^{(1)}[(\alpha,3),C_0],\mathbf{M}^{(1)}[(\alpha',1),C_0],\mathbf{M}^{(1)}[(\alpha',2),C_0],\mathbf{M}^{(1)}[(\alpha',3),C_0]$
restricted to the $6$ rows indexed by $\{(\alpha,1),\ldots,(\alpha',3)\}$. By
\eqref{eq:M1-eval} this column equals
\begin{equation}\label{eq:6-column}
\begin{bmatrix}\lambda_{\alpha\beta\gamma\delta}A^{(\alpha)}\\ \lambda_{\alpha'\beta\gamma\delta}A^{(\alpha')}\end{bmatrix}\cdot m(\beta,\gamma,\delta;j_0,k_0,\ell_0)\;\in\;\mathbb{R}^6.
\end{equation}
As $(j_0,k_0,\ell_0)$ varies over $[3]^3$, the vectors $m(\beta,\gamma,\delta;j_0,k_0,\ell_0)$
form a set of $27$ vectors in $\mathbb{R}^4$. By the proof of Lemma~\ref{lem:three-rows},
this set spans $\mathbb{R}^4$ (genericity on $A$). Hence the $3\times 4$ matrix
$\bigl[\lambda_{\alpha\beta\gamma\delta}A^{(\alpha)};\lambda_{\alpha'\beta\gamma\delta}A^{(\alpha')}\bigr]\in\mathbb{R}^{6\times 4}$
maps a spanning set of $\mathbb{R}^4$ to the columns of $\mathbf{M}^{(1)}$ restricted to the
$6$ rows above. Therefore the column span of this $6\times 4$ matrix coincides with the
column span of the $6\times 27$ block of $\mathbf{M}^{(1)}$ at columns indexed by
$\{((\beta,j_0),(\gamma,k_0),(\delta,\ell_0)):(j_0,k_0,\ell_0)\in[3]^3\}$.

We claim the $6\times 4$ matrix
\[
B(\beta,\gamma,\delta):=\begin{bmatrix}\lambda_{\alpha\beta\gamma\delta}A^{(\alpha)}\\ \lambda_{\alpha'\beta\gamma\delta}A^{(\alpha')}\end{bmatrix}
\]
has rank exactly $4$ if and only if $\lambda_{\alpha\beta\gamma\delta},\lambda_{\alpha'\beta\gamma\delta}\neq 0$.
Indeed, by (G3) (with $A^{(\alpha)},A^{(\alpha')}$ playing the roles), the union of the rows
of $A^{(\alpha)}$ and $A^{(\alpha')}$ spans $\mathbb{R}^4$ (six rows in $\mathbb{R}^4$, with
the first three forming a basis of a $3$-dim subspace and the bottom three spanning another
$3$-dim subspace; their union has dimension $4$ by (G3) applied with three distinct rows
chosen from these six, of which at most two come from one $A^{(\alpha)}$). Scaling the top
three rows by $\lambda_{\alpha\beta\gamma\delta}\neq 0$ and the bottom three by
$\lambda_{\alpha'\beta\gamma\delta}\neq 0$ does not change the rank.

Now the rank-$4$ constraint on $\mathbf{M}^{(1)}$ implies that the column span of the
$6\times 27$ block above sits in $V_1\cap\mathbb{R}^6_{\text{rows}\,\{(\alpha,i),(\alpha',i')\}}$,
which is again $4$-dimensional. The map $\mathbb{R}^4\to\mathbb{R}^6$ given by left
multiplication by $B(\beta,\gamma,\delta)$ has image equal to this $4$-dim space. Repeating
the construction with a different triple $(\beta',\gamma',\delta')$ yields the matrix
$B(\beta',\gamma',\delta')$ whose image is also a $4$-dim subspace of $\mathbb{R}^6$ — and
\emph{the same} subspace, because both equal the projection of $V_1$ onto the $6$ rows
$\{(\alpha,i),(\alpha',i')\}$ (this projection has dimension at least $4$ by
Lemma~\ref{lem:two-alpha-span} and at most $4$ by Lemma~\ref{lem:row-span}; hence exactly
$4$ and intrinsic to $\alpha,\alpha'$).

Therefore the column space of $B(\beta,\gamma,\delta)$ equals the column space of
$B(\beta',\gamma',\delta')$. Two $6\times 4$ matrices have the same column space iff one is
obtained from the other by right multiplication by a $4\times 4$ invertible matrix; but in
our setting both matrices have the special block form
$\bigl[\lambda_1 A^{(\alpha)};\lambda_2 A^{(\alpha')}\bigr]$. The relevant comparison is
between the columns $A^{(\alpha)}(:,p)$ and $\lambda_{\alpha'\beta\gamma\delta}/\lambda_{\alpha\beta\gamma\delta}\cdot A^{(\alpha')}(:,p)$.
Equality of column spans reads:
\[
\bigl[\lambda_{\alpha\beta\gamma\delta}A^{(\alpha)};\lambda_{\alpha'\beta\gamma\delta}A^{(\alpha')}\bigr]\cdot \mathbb{R}^4
=\bigl[\lambda_{\alpha\beta'\gamma'\delta'}A^{(\alpha)};\lambda_{\alpha'\beta'\gamma'\delta'}A^{(\alpha')}\bigr]\cdot \mathbb{R}^4.
\]
Equivalently, the linear map $\mathbb{R}^4\to\mathbb{R}^6$ has the same image. Multiplying
the top three coordinates by $1/\lambda_{\alpha\beta\gamma\delta}$ and bottom three by
$1/\lambda_{\alpha'\beta\gamma\delta}$ takes the first matrix to
$[A^{(\alpha)};A^{(\alpha')}]$. Doing the same to the second gives the same matrix only if
the diagonal scalings agree, i.e., the operator
$\mathrm{diag}(\lambda_{\alpha\beta\gamma\delta}/\lambda_{\alpha\beta'\gamma'\delta'},\ldots,
\lambda_{\alpha'\beta\gamma\delta}/\lambda_{\alpha'\beta'\gamma'\delta'},\ldots)$
preserves the column space of $[A^{(\alpha)};A^{(\alpha')}]$. By generic position of
$A^{(\alpha)},A^{(\alpha')}$ (specifically (G2) extended to ``no nontrivial diagonal scaling
of $[A^{(\alpha)};A^{(\alpha')}]$ preserves its column space''), this forces
\[
\frac{\lambda_{\alpha\beta\gamma\delta}}{\lambda_{\alpha\beta'\gamma'\delta'}}
=\frac{\lambda_{\alpha'\beta\gamma\delta}}{\lambda_{\alpha'\beta'\gamma'\delta'}}.
\]
This rearranges to $\lambda_{\alpha\beta\gamma\delta}/\lambda_{\alpha'\beta\gamma\delta}
=\lambda_{\alpha\beta'\gamma'\delta'}/\lambda_{\alpha'\beta'\gamma'\delta'}$, exactly
\eqref{eq:ratio}.
\end{proof}

\begin{lemma}[Mode-$1$ rank-$4$ implies rank-$1$ structure in $\alpha$]\label{lem:rank1-alpha}
Under the assumptions of Lemma~\ref{lem:lambda-firstindex}, there exists a function
$u:[n]\to\mathbb{R}^*$ and a function $\Lambda^{(1)}:[n]^3\to\mathbb{R}$ such that for every
quadruple $(\alpha,\beta,\gamma,\delta)$ with all four indices pairwise distinct,
\begin{equation}\label{eq:rank1-alpha}
\lambda_{\alpha\beta\gamma\delta}\;=\;u_\alpha\cdot\Lambda^{(1)}(\beta,\gamma,\delta).
\end{equation}
\end{lemma}

\begin{proof}
Fix a reference $\alpha_0\in[n]$. By Lemma~\ref{lem:lambda-firstindex}, for every
$\alpha\in[n]\setminus\{\alpha_0\}$ the ratio
$r_{\alpha,\alpha_0}=\lambda_{\alpha\beta\gamma\delta}/\lambda_{\alpha_0\beta\gamma\delta}$
is independent of any choice of $(\beta,\gamma,\delta)$ avoiding $\{\alpha,\alpha_0\}$ and
having distinct entries. Define $u_\alpha:=r_{\alpha,\alpha_0}\in\mathbb{R}^*$ for
$\alpha\neq\alpha_0$, and $u_{\alpha_0}:=1$. Define
$\Lambda^{(1)}(\beta,\gamma,\delta):=\lambda_{\alpha_0\beta\gamma\delta}$ for distinct
$\beta,\gamma,\delta\in[n]\setminus\{\alpha_0\}$ (and arbitrarily for the remaining
triples in a way consistent with the next steps; the values on these tuples will be
determined by Sections~\ref{sec:reverse-other-modes}--\ref{sec:reverse-finish}).

Then \eqref{eq:rank1-alpha} holds for $\alpha\neq\alpha_0$ by definition of $r_{\alpha,\alpha_0}$,
and for $\alpha=\alpha_0$ trivially.

It remains to handle quadruples $(\alpha,\beta,\gamma,\delta)$ where some of
$\{\alpha,\beta,\gamma,\delta\}$ coincide but not all four are equal. By assumption
$\lambda_{\alpha\beta\gamma\delta}\neq 0$ only when not all four are equal. We partition the
non-all-equal quadruples by their pattern of repetitions and address each case.

\textbf{Case~I:}\, all four distinct. \emph{Already handled above.}

\textbf{Case~II:}\, exactly one pair coincides among $\{\alpha,\beta,\gamma,\delta\}$. The
remaining cases (three or four coincidences) follow by similar reductions; we treat
Case~II in detail and remark on the rest.

Suppose $\alpha=\beta\neq\gamma\neq\delta\neq\alpha$, so the quadruple is
$(\alpha,\alpha,\gamma,\delta)$. To extract a relation, fix $\alpha_1\in[n]\setminus\{\alpha,\gamma,\delta\}$
(possible since $n\geq 5$). Take the column $C^*=\bigl((\alpha,j),(\gamma,k),(\delta,\ell)\bigr)$;
then \eqref{eq:M1-eval} reads
\[
\mathbf{M}^{(1)}\bigl[(\alpha,i),C^*\bigr]
=\lambda_{\alpha\alpha\gamma\delta}\,\bigl\langle A^{(\alpha)}(i,:),m^*\bigr\rangle,
\]
where $m^*$ is the vector of cofactor minors of
$[A^{(\alpha)}(j,:);A^{(\gamma)}(k,:);A^{(\delta)}(\ell,:)]$. Now consider the same
column $C^*$ and a different row $(\alpha_1,i)$:
\[
\mathbf{M}^{(1)}\bigl[(\alpha_1,i),C^*\bigr]
=\lambda_{\alpha_1\alpha\gamma\delta}\,\bigl\langle A^{(\alpha_1)}(i,:),m^*\bigr\rangle.
\]
Both expressions are projections of the row spaces in $V_1$. By
Lemma~\ref{lem:two-alpha-span} applied to $\alpha,\alpha_1$, the rows
$\{R^{(\alpha,i)},R^{(\alpha_1,i)}\}_{i}$ span $V_1$. Repeating the argument of
Lemma~\ref{lem:lambda-firstindex} with the column $C^*$ in place of $C_0$ — which is
permitted because by (G3) the matrix
$[A^{(\alpha)}(j,:);A^{(\gamma)}(k,:);A^{(\delta)}(\ell,:)]\in\mathbb{R}^{3\times 4}$ has
rank $3$ for any distinct $\alpha,\gamma,\delta$, so the cofactor vector $m^*\in\mathbb{R}^4$
is nonzero, and varying $(j,k,\ell)$ produces a spanning set of $\mathbb{R}^4$ — yields
\[
\frac{\lambda_{\alpha\alpha\gamma\delta}}{\lambda_{\alpha_1\alpha\gamma\delta}}
=\frac{\lambda_{\alpha\beta'\gamma'\delta'}}{\lambda_{\alpha_1\beta'\gamma'\delta'}}
=\frac{u_\alpha}{u_{\alpha_1}}=u_\alpha
\]
(taking $\alpha_1=\alpha_0$, so $u_{\alpha_1}=1$). By the analogous argument applied to the
mode-$2$ unfolding $\mathbf{M}^{(2)}$ — which has $\mathrm{rank}(\mathbf{M}^{(2)})\leq 4$ by
the hypothesis $\mathbf{F}(\mathbf{P})=0$ — we obtain a function $v:[n]\to\mathbb{R}^*$ with
$\lambda_{\alpha_1\alpha\gamma\delta}=v_\alpha\cdot\Lambda^{(2)}(\alpha_1,\gamma,\delta)$ for
$\alpha_1,\alpha,\gamma,\delta$ all distinct. Hence
\[
\lambda_{\alpha\alpha\gamma\delta}=u_\alpha\cdot\lambda_{\alpha_0\alpha\gamma\delta}
=u_\alpha\cdot v_\alpha\cdot\Lambda^{(2)}(\alpha_0,\gamma,\delta)
=u_\alpha\cdot\Lambda^{(1)}_{II}(\alpha,\gamma,\delta)
\]
where $\Lambda^{(1)}_{II}(\alpha,\gamma,\delta):=v_\alpha\cdot\Lambda^{(2)}(\alpha_0,\gamma,\delta)$
absorbs the dependence on $\alpha$ via $v_\alpha$.

\textbf{Case~III}\, (two coincident pairs, e.g., $\alpha=\beta$, $\gamma=\delta$, $\alpha\neq\gamma$):
the same argument applies; pick reference index $\alpha_0$ outside $\{\alpha,\gamma\}$ and
extract via mode-$1$ rank constraint the ratio $r_{\alpha,\alpha_0}=u_\alpha$.

\textbf{Case~IV}\, (three equal, one different; e.g., $\alpha=\beta=\gamma\neq\delta$): pick
reference $\alpha_0\notin\{\alpha,\delta\}$ and use the mode-$1$ rank constraint on the row
$(\alpha,i)$ with column $\bigl((\alpha,j),(\alpha,k),(\delta,\ell)\bigr)$. By
Lemma~\ref{lem:Q-skew} applied to the submatrix
$[A^{(\alpha)}(i,:);A^{(\alpha)}(j,:);A^{(\alpha)}(k,:);A^{(\delta)}(\ell,:)]$, two rows from
the same $A^{(\alpha)}$ may coincide (for the same first index repetition we chose), but
since $i,j,k\in[3]$ are not necessarily equal, we still have a $4\times 4$ matrix; by (G2)
applied to $A^{(\alpha)},A^{(\delta)}$ this matrix has rank $3$ unless the three top rows
are linearly dependent — they are by (G1) only when two of $i,j,k$ coincide. For columns
with $i,j,k$ all distinct in $[3]$, the cofactor vector $m^{**}$ is again nonzero and the
ratio extraction proceeds identically.

In all cases the conclusion is that
$\lambda_{\alpha\beta\gamma\delta}=u_\alpha\cdot\Lambda^{(1)}(\beta,\gamma,\delta)$ for some
function $\Lambda^{(1)}$ on the relevant set of triples.
\end{proof}

\subsection{Modes $2,3,4$ and joint rank-$1$ structure}\label{sec:reverse-other-modes}

\begin{lemma}[Rank-$1$ structures in each mode]\label{lem:rank1-each-mode}
Under the genericity hypothesis and assuming $\mathrm{rank}(\mathbf{M}^{(s)})\leq 4$ for
$s=1,2,3,4$, there exist functions $u,v,w,x:[n]\to\mathbb{R}^*$ and functions
$\Lambda^{(1)},\Lambda^{(2)},\Lambda^{(3)},\Lambda^{(4)}$ such that for every quadruple
$(\alpha,\beta,\gamma,\delta)$ with all four entries distinct,
\begin{align}
\lambda_{\alpha\beta\gamma\delta}&=u_\alpha\cdot\Lambda^{(1)}(\beta,\gamma,\delta),\label{eq:r1-alpha}\\
\lambda_{\alpha\beta\gamma\delta}&=v_\beta\cdot\Lambda^{(2)}(\alpha,\gamma,\delta),\label{eq:r1-beta}\\
\lambda_{\alpha\beta\gamma\delta}&=w_\gamma\cdot\Lambda^{(3)}(\alpha,\beta,\delta),\label{eq:r1-gamma}\\
\lambda_{\alpha\beta\gamma\delta}&=x_\delta\cdot\Lambda^{(4)}(\alpha,\beta,\gamma).\label{eq:r1-delta}
\end{align}
\end{lemma}

\begin{proof}
Apply Lemma~\ref{lem:rank1-alpha} to each mode $s\in\{1,2,3,4\}$ in turn. Mode $s$ uses the
unfolding $\mathbf{M}^{(s)}$ whose rows are indexed by $(\alpha_s,i_s)\in[n]\times[3]$, where
$\alpha_s$ is the $s$-th index. The argument is symmetric in modes; for mode $2$, one
expands the determinant defining $Q^{(\alpha\beta\gamma\delta)}_{ijk\ell}$ along the second
row instead of the first; for modes $3,4$ along the third and fourth rows.
\end{proof}

\begin{lemma}[Joint multiplicative form on quadruples with all distinct entries]\label{lem:jointly-multiplicative}
Under the assumptions of Lemma~\ref{lem:rank1-each-mode}, for every
$(\alpha,\beta,\gamma,\delta)\in[n]^4$ with all four entries distinct,
\begin{equation}\label{eq:full-multiplicative}
\lambda_{\alpha\beta\gamma\delta}=u_\alpha v_\beta w_\gamma x_\delta\cdot c
\end{equation}
for a single constant $c\in\mathbb{R}^*$ depending neither on $(\alpha,\beta,\gamma,\delta)$
nor on the matrices $A^{(\cdot)}$.
\end{lemma}

\begin{proof}
By \eqref{eq:r1-alpha} and \eqref{eq:r1-beta},
\[
u_\alpha\cdot\Lambda^{(1)}(\beta,\gamma,\delta)
=\lambda_{\alpha\beta\gamma\delta}
=v_\beta\cdot\Lambda^{(2)}(\alpha,\gamma,\delta).
\]
Hence $\Lambda^{(1)}(\beta,\gamma,\delta)/v_\beta=\Lambda^{(2)}(\alpha,\gamma,\delta)/u_\alpha$.
The left-hand side is independent of $\alpha$, the right-hand side independent of $\beta$, so
both equal a function $\Lambda^{(12)}(\gamma,\delta)$. That is,
$\Lambda^{(1)}(\beta,\gamma,\delta)=v_\beta\cdot\Lambda^{(12)}(\gamma,\delta)$, giving
\[
\lambda_{\alpha\beta\gamma\delta}=u_\alpha v_\beta\cdot\Lambda^{(12)}(\gamma,\delta).
\]
Substitute this into \eqref{eq:r1-gamma}: $u_\alpha v_\beta\Lambda^{(12)}(\gamma,\delta)=w_\gamma\Lambda^{(3)}(\alpha,\beta,\delta)$,
i.e., $\Lambda^{(12)}(\gamma,\delta)/w_\gamma=\Lambda^{(3)}(\alpha,\beta,\delta)/(u_\alpha v_\beta)$
which depends only on $\delta$. So $\Lambda^{(12)}(\gamma,\delta)=w_\gamma\Lambda^{(123)}(\delta)$,
giving $\lambda_{\alpha\beta\gamma\delta}=u_\alpha v_\beta w_\gamma\Lambda^{(123)}(\delta)$.
Substituting into \eqref{eq:r1-delta} gives $\Lambda^{(123)}(\delta)=x_\delta\cdot c$ for a
constant $c$, yielding \eqref{eq:full-multiplicative}.
\end{proof}

By rescaling $u_\alpha\mapsto c\cdot u_\alpha$ (and absorbing $c$ into $u$), we may set
$c=1$ throughout. Hence:
\begin{equation}\label{eq:multiplicative-distinct}
\lambda_{\alpha\beta\gamma\delta}=u_\alpha v_\beta w_\gamma x_\delta\quad\text{for all }(\alpha,\beta,\gamma,\delta)\in[n]^4\text{ with all entries distinct}.
\end{equation}

\subsection{Extending to non-all-equal quadruples with repetitions}\label{sec:reverse-finish}

We must verify \eqref{eq:multiplicative-distinct} continues to hold for every non-all-equal
quadruple, including those with some (but not all) coincident entries.

\begin{lemma}[Extension to repetitions]\label{lem:extension}
Assume \eqref{eq:multiplicative-distinct} for all four-distinct quadruples. Under the
genericity and rank assumptions, the same identity holds for every non-all-equal quadruple.
\end{lemma}

\begin{proof}
Let $(\alpha,\beta,\gamma,\delta)$ be a non-all-equal quadruple with at least one repetition.
We argue based on the pattern of repetitions.

\textbf{Pattern $\alpha=\beta$, others distinct from each other and from $\alpha$}: The
quadruple is $(\alpha,\alpha,\gamma,\delta)$ with $\alpha\neq\gamma\neq\delta\neq\alpha$.
Pick $\alpha_1\in[n]\setminus\{\alpha,\gamma,\delta\}$ (exists since $n\geq 5\geq 4$).
Apply mode-$1$ rank-$\leq 4$ to the row $(\alpha,i)$ at column
$C=((\alpha,j),(\gamma,k),(\delta,\ell))$ and to the row $(\alpha_1,i)$ at the same column.
By Lemma~\ref{lem:rank1-alpha}, the ratio
\[
\frac{\lambda_{\alpha\alpha\gamma\delta}}{\lambda_{\alpha_1\alpha\gamma\delta}}=u_\alpha/u_{\alpha_1}.
\]
Since $(\alpha_1,\alpha,\gamma,\delta)$ has all four distinct entries,
$\lambda_{\alpha_1\alpha\gamma\delta}=u_{\alpha_1}v_\alpha w_\gamma x_\delta$ by
\eqref{eq:multiplicative-distinct}. Hence
\[
\lambda_{\alpha\alpha\gamma\delta}=\frac{u_\alpha}{u_{\alpha_1}}\cdot u_{\alpha_1}v_\alpha w_\gamma x_\delta
=u_\alpha v_\alpha w_\gamma x_\delta.
\]
This matches the desired multiplicative form.

\textbf{Pattern $\alpha=\gamma$, others distinct}: $(\alpha,\beta,\alpha,\delta)$ with all
indices except positions $1,3$ distinct. Pick $\alpha_1\in[n]\setminus\{\alpha,\beta,\delta\}$.
Apply mode-$1$ rank to compare rows $(\alpha,i),(\alpha_1,i)$ at column
$C=((\beta,j),(\alpha,k),(\delta,\ell))$:
\[
\frac{\lambda_{\alpha\beta\alpha\delta}}{\lambda_{\alpha_1\beta\alpha\delta}}=u_\alpha/u_{\alpha_1}.
\]
The denominator's quadruple $(\alpha_1,\beta,\alpha,\delta)$ has all four distinct, so it
equals $u_{\alpha_1}v_\beta w_\alpha x_\delta$ by \eqref{eq:multiplicative-distinct}. Hence
\[
\lambda_{\alpha\beta\alpha\delta}=u_\alpha v_\beta w_\alpha x_\delta.
\]

\textbf{Pattern $\alpha=\delta$}: analogous; exchange $\delta$ for $\alpha$ in the argument.

\textbf{Pattern $\beta=\gamma$, others distinct}: $(\alpha,\beta,\beta,\delta)$ with $\alpha,\beta,\delta$
distinct. Apply the mode-$2$ rank-$\leq 4$ argument: the analog of
Lemma~\ref{lem:rank1-alpha} for mode $2$ gives a ratio
$\lambda_{\alpha\beta\beta\delta}/\lambda_{\alpha\beta_1\beta\delta}=v_\beta/v_{\beta_1}$
with $\beta_1\in[n]\setminus\{\alpha,\beta,\delta\}$. The denominator's quadruple
$(\alpha,\beta_1,\beta,\delta)$ has all four distinct, so equals
$u_\alpha v_{\beta_1}w_\beta x_\delta$. Hence
$\lambda_{\alpha\beta\beta\delta}=u_\alpha v_\beta w_\beta x_\delta$.

\textbf{Pattern $\beta=\delta$}: analogous via mode $2$ or mode $4$.

\textbf{Pattern $\gamma=\delta$}: analogous via mode $3$ or mode $4$.

\textbf{Pattern with two pairs}, e.g., $\alpha=\beta$ and $\gamma=\delta$, $\alpha\neq\gamma$:
$(\alpha,\alpha,\gamma,\gamma)$. Pick $\alpha_1\notin\{\alpha,\gamma\}$ (which exists since
$n\geq 5\geq 3$). Apply mode-$1$:
$\lambda_{\alpha\alpha\gamma\gamma}/\lambda_{\alpha_1\alpha\gamma\gamma}=u_\alpha/u_{\alpha_1}$.
Now $\lambda_{\alpha_1\alpha\gamma\gamma}$ has the pattern $\beta=\delta$ which we just
handled: $\lambda_{\alpha_1\alpha\gamma\gamma}=u_{\alpha_1}v_\alpha w_\gamma x_\gamma$.
Hence $\lambda_{\alpha\alpha\gamma\gamma}=u_\alpha v_\alpha w_\gamma x_\gamma$.

\textbf{Pattern $\alpha=\gamma$, $\beta=\delta$}: $(\alpha,\beta,\alpha,\beta)$ with
$\alpha\neq\beta$. Pick $\alpha_1\notin\{\alpha,\beta\}$. By mode-$1$,
$\lambda_{\alpha\beta\alpha\beta}/\lambda_{\alpha_1\beta\alpha\beta}=u_\alpha/u_{\alpha_1}$.
$\lambda_{\alpha_1\beta\alpha\beta}$ has pattern $\beta=\delta$ (positions $2,4$ equal): it
equals $u_{\alpha_1}v_\beta w_\alpha x_\beta$ by the previous case. Hence
$\lambda_{\alpha\beta\alpha\beta}=u_\alpha v_\beta w_\alpha x_\beta$.

\textbf{Pattern $\alpha=\delta$, $\beta=\gamma$}: $(\alpha,\beta,\beta,\alpha)$ with
$\alpha\neq\beta$: analogous, equals $u_\alpha v_\beta w_\beta x_\alpha$.

\textbf{Pattern with three equal, fourth different}, e.g., $\alpha=\beta=\gamma\neq\delta$:
$(\alpha,\alpha,\alpha,\delta)$. Pick $\alpha_1\notin\{\alpha,\delta\}$. By mode-$1$,
$\lambda_{\alpha\alpha\alpha\delta}/\lambda_{\alpha_1\alpha\alpha\delta}=u_\alpha/u_{\alpha_1}$.
$\lambda_{\alpha_1\alpha\alpha\delta}$ has pattern $\beta=\gamma$: equals
$u_{\alpha_1}v_\alpha w_\alpha x_\delta$ by an earlier case. Hence
$\lambda_{\alpha\alpha\alpha\delta}=u_\alpha v_\alpha w_\alpha x_\delta$. Other patterns
with three equal entries are handled symmetrically.

In every pattern of repetitions (excluding the all-equal case, which is excluded by
hypothesis), $\lambda_{\alpha\beta\gamma\delta}=u_\alpha v_\beta w_\gamma x_\delta$.
\end{proof}

\begin{proposition}[Reverse direction of (P3)]\label{prop:rev}
Suppose $\mathbf{F}(\mathbf{P})=0$ and the genericity hypothesis on $A^{(1)},\ldots,A^{(n)}$
holds. Then there exist $u,v,w,x\in(\mathbb{R}^*)^n$ such that for every non-all-equal
quadruple $(\alpha,\beta,\gamma,\delta)$,
$\lambda_{\alpha\beta\gamma\delta}=u_\alpha v_\beta w_\gamma x_\delta$.
\end{proposition}

\begin{proof}
By assumption, every $5\times 5$ minor of $\mathbf{M}^{(s)}$ vanishes for $s=1,2,3,4$. Hence
$\mathrm{rank}(\mathbf{M}^{(s)})\leq 4$. By Lemma~\ref{lem:rank1-each-mode}, the array
$\lambda$ has rank-$1$ structure in each mode. By
Lemma~\ref{lem:jointly-multiplicative} with the constant absorbed into $u$,
$\lambda_{\alpha\beta\gamma\delta}=u_\alpha v_\beta w_\gamma x_\delta$ for all
four-distinct quadruples. By Lemma~\ref{lem:extension}, the same multiplicative form holds
for every non-all-equal quadruple, including those with repetitions.

Finally, $u_\alpha,v_\beta,w_\gamma,x_\delta\in\mathbb{R}^*$ because the values
$\lambda_{\alpha\beta\gamma\delta}$ are nonzero on quadruples with four distinct entries
(since $n\geq 5$, four-distinct quadruples exist for every $\alpha\in[n]$); hence each
factor must be nonzero.
\end{proof}

\section{Conclusion}\label{sec:conclusion}

\begin{theorem}\label{thm:main}
For every $n\geq 5$ and Zariski-generic $A^{(1)},\ldots,A^{(n)}\in\mathbb{R}^{3\times 4}$,
the polynomial map $\mathbf{F}$ from Definition~\ref{def:F} satisfies properties (P1),
(P2), (P3). The number of coordinate functions is $N=4\binom{3n}{5}\binom{(3n)^3}{5}$, each
of degree exactly $5$ with $\pm 1$ integer coefficients.
\end{theorem}

\begin{proof}
Property (P1) and (P2) are Proposition~\ref{prop:P1P2}. Property (P3) is the conjunction of
Proposition~\ref{prop:fwd} (forward direction) and Proposition~\ref{prop:rev} (reverse
direction).
\end{proof}

\begin{remark}
The construction is essentially the rank-$\leq 4$ flattening test for the $4$-tensor
$\mathbf{P}$ of size $3n$ in each mode. The rank bound of $4$ matches the column dimension
of the $A^{(\alpha)}\in\mathbb{R}^{3\times 4}$, reflecting that each mode-$s$ row of
$\mathbf{P}$ lies in the $4$-dim row span generated by cofactor expansion of the underlying
$4\times 4$ determinants. The constant $5$ in ``$5\times 5$ minors'' is therefore optimal:
$4\times 4$ minors do not in general vanish, while $5\times 5$ minors vanish exactly on the
multiplicative locus described by the rank-$4$ flattening conditions.
\end{remark}

\end{document}
